\chapter{Background}
\label{ch:background}

Large language models routinely produce fluent answers that are not
supported by the facts of the task. Such failures are typically grouped
under the broad label of ``hallucination''. The term covers several
distinct mechanisms: unfaithfulness to source material, unsupported
factual claims, gaps in the model's knowledge, faulty reasoning, and
answers encouraged by benchmarks that reward guessing over
admitting ignorance~\cite{ji_hallucination_survey_2023,%
huang_llm_hallucination_survey_2023,%
maynez_faithfulness_factuality_2020,%
kalai_why_language_models_hallucinate_2025,%
augenstein_factuality_challenges_2024}. This thesis uses the word as a
convenient cover term, but does not try to detect every kind. The target is
narrower: answers that are unreliable because the model is uncertain about
what it just said.

\section{Hallucinations and confabulations}

The semantic entropy paper~\cite{semantic_entropy_nature_2024}, on which
this thesis is based, sharpens the target to \emph{confabulations}:
fluent answers that are both incorrect and arbitrary, where ``arbitrary''
means that the answer is sensitive to irrelevant variation such as the
random seed used during decoding. In addition to confabulations, several
other types of hallucinations can occur: a model may be consistently
wrong because its training data contains a misconception, because the
problem requires reasoning it cannot perform, or because evaluation
incentives discourage abstention. Semantic entropy does not address
these cases; it specifically detects instability across sampled
meanings.

This narrower framing motivates the central hypothesis of the thesis,
again taken from Farquhar et
al.~\cite{semantic_entropy_nature_2024}: if the model knows
the answer, stochastic samples should vary mostly in wording while
preserving meaning. If the model is confabulating, the samples should
split across several incompatible meanings. Measuring uncertainty over
meanings rather than over surface strings should therefore be a more
informative signal for free-form question answering, following
semantic uncertainty and semantic entropy
work~\cite{kuhn_semantic_uncertainty_2023,semantic_entropy_nature_2024}.
In particular, uncertainty-based detection is, by construction, unable
to flag high-certainty hallucinations or systematic
errors~\cite{simhi_choke_2025,ma_semantic_energy_2025}.

\section{Uncertainty in free-form generation}

When a model has to pick from a fixed, short list of answers, say
``cat'', ``dog'', or ``horse'', measuring its uncertainty is easy: it
assigns a percentage to each option, and those percentages already say
how sure it is. A language model has no such tidy list. It instead
produces whole sentences, and the catch is that the \emph{same fact} can
be written in many ways (``Paris'', ``Paris, France'', ``the capital is
Paris''), while two sentences that look almost identical can mean
different things. So counting how often the exact words differ is a poor
guide to whether the model is really unsure. Split the model's probability
mass across ten paraphrases of one fact and the token-level numbers look
uncertain, even though the model has said the same thing ten times. Sequence-level entropy or
token-level variation can therefore be misleading for natural-language
answers, as discussed in work on uncertainty in natural language
generation~\cite{baan_uncertainty_nlg_2023} and semantic
uncertainty~\cite{kuhn_semantic_uncertainty_2023}.
The thesis distinguishes three levels of variation accordingly:

\begin{enumerate}
  \item \emph{lexical variation}: different token sequences or
    paraphrases (e.g.\ ``Paris'' vs.\ ``Paris, France'');
  \item \emph{semantic variation}: different meanings or propositions
    (e.g.\ ``Paris'' vs.\ ``Lyon'');
  \item \emph{factual correctness}: whether the chosen meaning is
    correct for the task.
\end{enumerate}

Semantic entropy operates at the second level. It estimates uncertainty
over semantic answers, not over truth directly; correctness evaluation
is a separate step that compares the selected answer against a reference
or a judge in semantic entropy~\cite{semantic_entropy_nature_2024}.
Recent surveys~\cite{shorinwa_uq_llm_survey_2024,liu_uq_confidence_calibration_survey_2025}
place
this family of methods within uncertainty quantification for LLMs,
alongside calibration, verbalized confidence, internal-state methods,
and other sampling-based estimators, and systematic benchmarks now
compare many such estimators head to
head~\cite{lm_polygraph_vashurin_2025}.

\section{Confidence, calibration, and abstention}

The motivation for estimating uncertainty is not only to attach a score
to an answer, but to decide what should happen next: answer, abstain,
defer to retrieval, request human review, or invoke a more expensive
verification procedure. A related notion is \emph{calibration}. A model is well calibrated when its
stated confidence matches how often it is actually right: among all answers
given with 90\,\% confidence, roughly 90\,\% should be correct. Calibration
and accuracy are different properties, and a model can have one without the
other. The view taken here needs neither. Selective prediction only requires
that risky answers be \emph{ranked} above safe ones. This
makes selective prediction and
rejection-based evaluation the natural evaluation lens for the thesis.
Accuracy alone is insufficient, because many benchmarks reward a lucky
guess more than an honest admission of uncertainty, a problem discussed
by Kalai et al.~\cite{kalai_why_language_models_hallucinate_2025}.
Rejection-accuracy curves and their AURAC summary are better aligned
with the practical question of whether uncertain answers can be
filtered out. Related work has also
examined whether semantic entropy can be used as a training signal to
encourage abstention: Tjandra et
al.~\cite{tjandra_abstain_semantic_entropy_2024} fine-tune models on it
directly. That is a different use of the same quantity, and it is not what
this thesis does. Here the score is only ever computed on a finished answer
from a model that is left untouched; no gradient is taken through it.

A number of related works have approached the same problem without
explicit semantic clustering. P(True)~\cite{kadavath_language_models_know_2022}
asks the model to judge whether
its own answer is correct, probing verbalized self-knowledge rather than
sample diversity. SelfCheckGPT~\cite{manakul_selfcheckgpt_2023}
instead samples multiple responses and checks them for consistency
without relying on external databases or access to token
probabilities. Semantic entropy is
closest in spirit to these sampling-based methods, but, unlike them, it
explicitly aggregates samples into semantic clusters before computing
an entropy score~\cite{kuhn_semantic_uncertainty_2023,semantic_entropy_nature_2024}.
These methods span a spectrum, from purely internal (P(True) reads the
model's own verdict) to purely behavioral (SelfCheckGPT compares samples).
Semantic entropy sits between them. It adds one step none of the others
perform --- deciding when two answers mean the same thing --- and the value
of exactly that step is what this study measures.

A third family reads the model's \emph{hidden states} rather than its words.
The observation behind it is that a model's internal activations often carry
information about the reliability of an answer that its output text does not
express. Burns et al.~\cite{burns_latent_knowledge_2023} recover a
truth-like direction in activation space without any labels; Azaria and
Mitchell~\cite{azaria_internal_state_lying_2023} train a small classifier on
hidden states to tell true from false statements; and
INSIDE~\cite{chen_inside_2024} scores hallucinations from the geometry of
sentence embeddings taken from inside the model. The practical appeal is cost:
these methods need one forward pass, not $M$ samples and $M^2$ entailment
comparisons. Semantic Entropy Probes~\cite{kossen_semantic_entropy_probes_2024}
join the two families by training exactly such a probe to predict semantic
entropy itself, so that the expensive signal is paid for once at training time
and read cheaply thereafter. Both a probe of this kind and a
correctness-trained variant are evaluated in
Chapter~\ref{ch:results} as single-pass comparators; they are defined in
Sections~\ref{sec:sep-method} and~\ref{sec:other-estimators}. Note that
training such a probe is a separate, supervised step layered on top of a frozen
generator, and it does not change the generator either.

\section{Semantic equivalence and entailment}

For computing the semantic entropy, the key operational step is
deciding when two sampled answers should be treated as the same answer.
Semantic entropy adopts semantic equivalence rather than exact string
equality: in Farquhar et al.'s original
formulation~\cite{semantic_entropy_nature_2024}, two answers belong to
the same group when they entail one another in the context of the question.
This bidirectional-entailment
criterion is a practical approximation to equality of meaning.

Natural Language Inference (NLI) is the standard tool for this. An NLI
model reads two sentences and labels their relationship as
\emph{entailment} (the first implies the second), \emph{contradiction},
or \emph{neutral}. Checking entailment in both directions is then a
practical stand-in for ``these two answers mean the same thing.'' The
specific NLI models used here are DeBERTa-v3 \emph{cross-encoders} (in
two sizes, base and large) fine-tuned on the standard SNLI and MultiNLI inference
datasets~\cite{sentence_transformers_nli_deberta_v3_base_2022,%
he_deberta_2021,he_debertav3_2021,williams_multinli_2018}. ``Cross-encoder''
just means the two sentences are fed into the model
\emph{together}, so it can directly compare them; the alternative, a
\emph{bi-encoder}, would encode each sentence on its own and compare the
results afterwards. The cross-encoder is more accurate for this kind of
pairwise judgement but slower, because every pair of answers has to be
run through the model.
The exact bidirectional-entailment rule and the recorded backend choices
are settled in Chapters~\ref{ch:semantic_entropy} and~\ref{ch:results}.

\section{Research gap}

Prior evaluations of semantic entropy report strong aggregate AUROC but
leave three questions unresolved. First, the original evaluation compares
semantic entropy against naive entropy over raw sampled strings, but not
against a normalized surface-form
baseline~\cite{semantic_entropy_nature_2024}, and the survey and
benchmark literature that follows it inherits that set of
comparisons~\cite{shorinwa_uq_llm_survey_2024,%
liu_uq_confidence_calibration_survey_2025,lm_polygraph_vashurin_2025}. It
is therefore unclear how much of the
reported gain over ``naive'' entropy is attributable to semantic
clustering rather than to simple answer normalization. Second, short-answer and long-form results are usually reported together
rather than as a controlled contrast. Yet these are exactly the two regimes
in which surface-level and meaning-level uncertainty should come apart.
Recent work does report that semantic entropy weakens as one-sentence outputs
grow longer~\cite{nguyen_beyond_semantic_entropy_2025}, but it does not
isolate the short-versus-long contrast under one protocol. Third,
the dependence of sampling-based detection on the decoding temperature,
the very parameter that controls sample diversity, has been
examined only in isolated ablations (including in the same
work~\cite{nguyen_beyond_semantic_entropy_2025}) rather than
characterized systematically across tasks and detector families. This
thesis targets these three gaps directly.

Addressing them fixes the empirical scope to three short-answer tasks,
namely open-domain factual question answering with
TriviaQA~\cite{joshi_triviaqa_2017} and NQ-Open~\cite{lee_nq_open_2019}
and arithmetic word problems with SVAMP~\cite{patel_svamp_2021}, and
one long-form task, a FactScore-style biography
task~\cite{min_factscore_2023}. The dataset rationale, the datasets left
outside the recorded study, representative records, and the exact prompt
construction are given in Chapter~\ref{ch:experimental_design} and
Appendix~\ref{app:prompts}.

Taken together, the literature places semantic entropy between cheap
surface signals and expensive factual verification. Testing it fairly means
keeping generation, clustering, scoring and evaluation separate, so this
thesis builds a small, verifiable implementation first.
